% Appendix with a bibliography, lists of figures and tables, and a glossary. The appendix is lettered
% (A, B, ...), left out of the outline, and uses the standin frame style.
\documentclass[aspectratio=169,10pt]{beamer}
\usepackage[backend=biber, style=numeric]{biblatex}
\addbibresource{joby-appendix.bib}
\usepackage[style=list, nonumberlist]{glossaries}
\makeglossaries
\newglossaryentry{ral}{name={RAL}, description={Redundant Abstraction Layer}}
\newglossaryentry{fcc}{name={FCC}, description={Flight Control Computer}}
\usetheme{joby}
\title{Appendix test}
\author{Test Author}
\begin{document}
\maketitle

\begin{frame}[toc]{Outline}
	\tableofcontents
\end{frame}

\section{Main body}
\begin{frame}{Content}
	The \gls{ral} runs on the \gls{fcc}; see \cite{knuth1984,lamport1994}.
	\begin{figure}
		\rule{3cm}{1cm}
		\caption{A first figure}
	\end{figure}
	\begin{table}
		\begin{tabular}{cc}1 & 2\end{tabular}
		\caption{A first table}
	\end{table}
\end{frame}

\appendix
\section{Bibliography}
\begin{frame}[allowframebreaks]{References}
	\printbibliography[heading=none]
\end{frame}

\section{Lists of Figures and Tables}
\begin{frame}{List of figures}
	\listoffigures
\end{frame}
\begin{frame}{List of tables}
	\listoftables
\end{frame}

\section{Glossary}
\begin{frame}[allowframebreaks]{Glossary}
	\printglossaries
\end{frame}
\end{document}
